% ============================================================================
% main.tex - DIE QUELLE: IWT UND Q
% ============================================================================

\documentclass[11pt,a4paper]{book}

\usepackage[utf8]{inputenc}
\usepackage[ngerman]{babel}
\usepackage{amsmath, amssymb, amsfonts}
\usepackage{graphicx}
\usepackage{hyperref}
\usepackage{geometry}
\usepackage{fancyhdr}
\usepackage{tocloft}
\usepackage{float}
\usepackage{caption}
\usepackage{subcaption}
\usepackage{booktabs}
\usepackage{array}
\usepackage{enumitem}
\usepackage{listings}
\usepackage{xcolor}

\geometry{a4paper, margin=2.5cm}

% ============================================================================
% FANCYHDR - KORREKTUR
% ============================================================================

\setlength{\headheight}{13.7pt}
\addtolength{\topmargin}{-0.1pt}

% ============================================================================

\hypersetup{
    colorlinks=true,
    linkcolor=blue,
    urlcolor=blue,
    citecolor=blue
}

\pagestyle{fancy}
\fancyhf{}
\fancyhead[L]{\leftmark}
\fancyhead[R]{\thepage}
\renewcommand{\headrulewidth}{0.4pt}

\setlength{\parindent}{0pt}
\setlength{\parskip}{1em}

% ============================================================================
% TITEL
% ============================================================================

\title{
    {\Huge \textbf{Die Quelle}}\\
    {\Large Information, Selbstorganisation und Q}\\
    \vspace{0.5cm}
    {\large Die IWT als fundamentale Ur-Theorie}
}
\author{Michael Czybor}
\date{27. Juli 2026}

\begin{document}

\maketitle
\tableofcontents

% ============================================================================
% EINLEITUNG
% ============================================================================

\chapter{Einleitung}
\label{ch:einleitung}

% ============================================================================
% kapitel/einleitung.tex
% ============================================================================

\section{Das Problem der modernen Physik}

Die moderne Physik steht vor einem Rätsel. Sie verfügt mit der Quantenmechanik und der Allgemeinen Relativitätstheorie über zwei der erfolgreichsten Theorien der Wissenschaftsgeschichte. Aber sie sind fundamental inkompatibel.

Die vorgeschlagenen Lösungen – Stringtheorie, Schleifenquantengravitation, dunkle Materie, dunkle Energie – haben das Problem nicht gelöst.

Es fehlt ein einheitliches Fundament.

\section{Die Lösung: IWT und Q}

Dieses Dokument stellt die Informations-Weber-Theorie (IWT) vor.

Sie hat nur eine fundamentale Größe: \textbf{Information} \( I \).

Aus der Information emergiert alles andere: Raum, Zeit, Materie, Energie, Gravitation, Quantenmechanik, Elektrodynamik.

Die organisierende Kraft ist \textbf{Q}: die Quelle der Selbstorganisation.

\section{Zwei Quellen}

Dieses Dokument ruht auf zwei Säulen:

\begin{enumerate}
    \item \textbf{Die Theorie:} Axiome, Herleitungen, mathematische Struktur
    \item \textbf{Die Simulation:} Numerische Bestätigung, Ergebnisse, Visualisierung
\end{enumerate}

Die Theorie sagt, was sein muss. Die Simulation zeigt, dass es ist.

\section{Die zentrale Aussage}

\textbf{Q ist die Quelle.}

Alles emergiert aus Q.

\begin{itemize}
    \item Die DBT emergiert aus Q (Bohm-Potential)
    \item Die WG emergiert aus Q (Gravitation)
    \item Die WED emergiert aus Q (Elektrodynamik)
\end{itemize}

Q ist das Prinzip der Selbstorganisation.

Q ist die Quelle der Physik.

% ============================================================================
% TEIL I: THEORIE
% ============================================================================

\part{Die Theorie}

\chapter{Die Informations-Weber-Theorie (IWT)}
\label{ch:iwt}

% ============================================================================
% kapitel/iwt.tex
% ============================================================================

\section{Die grundlegende Größe: Information}

Die IWT hat nur eine fundamentale Größe: die Information \( I_k^{(n)} \).

\[
I_k^{(n)} \in \mathbb{C}
\]

\begin{itemize}
    \item \( k \) = Index des Informationsknotens
    \item \( n \) = diskreter Zeitindex
    \item \( I_k^{(n)} \) = komplexer Informationswert
\end{itemize}

Die Amplitude \( |I| \) kodiert die Stärke der Information.

Die Phase \( \arg(I) \) kodiert die kohärenten Relationen.

\section{Die Informationserhaltung}

Die Gesamtinformation ist erhalten:

\[
\sum_{k=1}^{N} |I_k^{(n)}|^2 = \text{const}
\]

Das ist Axiom 2 der IWT.

\section{Die Kopplungsmatrix}

Die Kopplungsmatrix beschreibt die Wechselwirkung zwischen Knoten:

\[
K_{kl} = \frac{1}{d_{kl}^{3-D}}
\]

Dabei ist \( d_{kl} \) die fraktale Distanz und \( D \) die fraktale Dimension.

\section{Die Dynamik}

Die Dynamik folgt aus dem Variationsprinzip:

\[
I_k^{(n+1)} = I_k^{(n)} + T \cdot \Phi_k
\]

Dabei ist \( \Phi_k \) die Kraft, die aus dem Lagrange-Funktional folgt.

\section{Die Rolle von Q}

Q ist nicht Teil der Axiome.

Q ist ein Resultat der Selbstorganisation.

Q emergiert aus der Information.

\chapter{Q: Die Quelle der Organisation}
\label{ch:q}

% ============================================================================
% kapitel/q.tex
% ============================================================================

\section{Q: Die Quelle der Organisation}

Q ist das zentrale Prinzip der IWT.

Es ist nicht das Bohm-Potential.

Es ist die Organisation der Information.

\section{Die Definition von Q}

Q ist definiert als das Prinzip, das die Information so organisiert, dass sie stabil ist.

\[
Q = \text{Organisation}(I)
\]

Es ist die Quelle aller Struktur.

\section{Die Herleitung von Q}

Q folgt aus zwei Prinzipien:

\begin{enumerate}
    \item \textbf{Informationserhaltung:} \( \sum |I|^2 = \text{const} \)
    \item \textbf{Minimale Spannung:} Die Information wird so organisiert, dass die Änderung zwischen benachbarten Knoten minimal ist.
\end{enumerate}

Die Lagrange-Multiplikator-Methode ergibt:

\[
\Delta^2 I = \lambda \cdot I
\]

Daraus folgt Q als das Potential, das diese Bedingung erfüllt.

\section{Q als Quelle}

Q ist die Quelle von:

\begin{itemize}
    \item \textbf{DBT:} Q ist das Bohm-Potential
    \item \textbf{WG:} \( F = -\nabla Q \) (Gravitation)
    \item \textbf{WED:} \( F = -\nabla Q \) (Elektrodynamik)
\end{itemize}

Alles kommt aus Q.

\section{Q als Selbstorganisation}

Q ist das Prinzip der Selbstorganisation.

Die Information organisiert sich selbst.

Q ist der Name dieser Organisation.

Q ist die Quelle.

\chapter{Emergenz der Physik aus Q}
\label{ch:emergenz}

% ============================================================================
% kapitel/emergenz.tex
% ============================================================================

\section{Die Emergenz der Physik aus Q}

Aus Q emergiert die gesamte Physik.

\section{Die DBT (De-Broglie-Bohm-Theorie)}

Die DBT ist die Quantentheorie.

Ihr Kern ist das Bohm-Potential:

\[
Q = -\frac{\hbar^2}{2m} \frac{\nabla^2 \sqrt{\rho}}{\sqrt{\rho}}
\]

Das ist Q in der DBT.

Die DBT ist Q.

\section{Die WG (Weber-Gravitation)}

Die WG ist die Gravitation.

Ihre Kraft ist:

\[
F_{\text{WG}} = -\frac{GMm}{r^2} \left[1 - \frac{\dot{r}^2}{c^2} + \frac{1}{2}\frac{r\ddot{r}}{c^2}\right]
\]

Aus Q folgt:

\[
F_{\text{WG}} = -\nabla Q
\]

Die WG ist der Gradient von Q.

\section{Die WED (Weber-Elektrodynamik)}

Die WED ist die Elektrodynamik.

Ihre Kraft ist:

\[
F_{\text{WED}} = \frac{q_1 q_2}{4\pi\epsilon_0 r^2} \left[1 - \frac{\dot{r}^2}{c^2} + 2\frac{r\ddot{r}}{c^2}\right]
\]

Aus Q folgt:

\[
F_{\text{WED}} = -\nabla Q
\]

Die WED ist der Gradient von Q.

\section{Die Einheit}

Die gesamte Physik ist Q.

\begin{itemize}
    \item DBT = Q
    \item WG = \( -\nabla Q \)
    \item WED = \( -\nabla Q \)
\end{itemize}

Q ist die Quelle.

\chapter{Axiome der IWT}
\label{ch:axiome}

% ============================================================================
% kapitel/axiome.tex
% ============================================================================

\section{Die Axiome der IWT}

Die IWT hat neun Axiome.

\subsection{Axiom 1: Diskretes Informationsfeld}

Es existiert ein skalares, diskretes Informationsfeld \( I_k^{(n)} \).

\subsection{Axiom 2: Informationserhaltung}

\[
\sum_{k=1}^{N} |I_k^{(n)}|^2 = \text{const}
\]

\subsection{Axiom 3: Informationsmetrik}

Die Metrik emergiert aus der Kopplungsmatrix:

\[
g_{kl}^{(n)} = \frac{K_{kl}^{(n)}}{\sqrt{K_{kk}^{(n)} K_{ll}^{(n)}}}
\]

\subsection{Axiom 4: Variationsprinzip}

Die Dynamik folgt aus einem diskreten Lagrange-Funktional.

\subsection{Axiom 5: Dynamische Metrik}

Die Metrik entwickelt sich dynamisch.

\subsection{Axiom 6: Energie als emergenter Informationsfluss}

Energie ist eine Erhaltungsgröße des Informationsflusses.

\subsection{Axiom 7: Naturkonstanten als emergente Skalierungsparameter}

Die Naturkonstanten folgen aus der Informationsdynamik.

\subsection{Axiom 8: Emergenz von Raum, Zeit und Dynamik}

Raum, Zeit und Dynamik emergieren aus der Information.

\subsection{Axiom 9: Universelle Gültigkeit}

Die IWT gilt auf allen Skalen.

\section{Q ist kein Axiom}

Q ist nicht in den Axiomen enthalten.

Q ist ein Resultat der Selbstorganisation.

Q ist die Quelle.

Q emergiert aus der Information.

% ============================================================================
% TEIL II: SIMULATION
% ============================================================================

\part{Die Simulation}

\chapter{Die Simulation der IWT}
\label{ch:simulation}

% ============================================================================
% kapitel/simulation.tex
% ============================================================================

\section{Die Simulation der IWT}

Die IWT wurde numerisch implementiert.

Die Simulation verwendet:

\begin{itemize}
    \item \( N = 4096 \) Knoten
    \item \( 500 \) Iterationen
    \item GPU-Beschleunigung (OpenCL)
    \item Leapfrog-Integrator (symplektisch)
\end{itemize}

\section{Die eingefrorenen Funktionen}

Die Simulation hat zwei eingefrorene Funktionen:

\begin{enumerate}
    \item \textbf{Fluktuationen (Vakuum):} Nur im Vakuum aktiv
    \item \textbf{Energiesenke (Strukturen):} Nur in Strukturen aktiv
\end{enumerate}

Diese Funktionen sind stabil und werden nicht mehr geändert.

\section{Die Parameter}

Die Parameter der Simulation sind:

\begin{align*}
    \text{SCALE} &= 0.70710678 \\
    \text{ALPHA\_0} &= 1e-2 \\
    \text{RHO\_0} &= 1e-6 \\
    \text{RHO\_MIN} &= 1e-8
\end{align*}

\section{Die Stabilität}

Die Simulation ist stabil:

\[
\sum |I|^2 = \text{const}
\]

Axiom 2 ist erfüllt.

\chapter{Ergebnisse der Simulation}
\label{ch:ergebnisse}

% ============================================================================
% kapitel/ergebnisse.tex
% ============================================================================

\section{Die Ergebnisse der Simulation}

Die Simulation zeigt:

\subsection{Strukturbildung}

Es gibt Strukturen:

\[
I_{\text{max}} > I_{\text{min}}
\]

Die Strukturen sind stabil.

\subsection{Q ist negativ}

Der Mittelwert von Q ist negativ:

\[
\langle Q \rangle < 0
\]

Q wirkt bindend.

\subsection{Die Gravitation emergiert aus Q}

Die Gravitation ist der Gradient von Q:

\[
F = -\nabla Q
\]

Die Simulation zeigt: Q ist die Quelle der Gravitation.

\subsection{Die Einheit der Physik}

Die Simulation zeigt, dass Q die Quelle von allem ist:

\begin{itemize}
    \item DBT = Q
    \item WG = \( -\nabla Q \)
    \item WED = \( -\nabla Q \)
\end{itemize}

\section{Die Bestätigung}

Die Simulation bestätigt die Theorie.

Die Theorie sagt, was sein muss.

Die Simulation zeigt, dass es ist.

\chapter{Die Bestätigung der Theorie}
\label{ch:bestaetigung}

% ============================================================================
% kapitel/bestaetigung.tex
% ============================================================================

\section{Die Bestätigung der Theorie}

Die Simulation hat die Theorie bestätigt.

\subsection{Axiom 2 ist erfüllt}

\[
\sum |I|^2 = \text{const}
\]

Die Information bleibt erhalten.

\subsection{Q emergiert}

Q emergiert aus der Selbstorganisation der Information.

Q ist nicht postuliert. Es ist ein Resultat.

\subsection{Die Gravitation emergiert aus Q}

Die Gravitation ist der Gradient von Q.

Die Simulation zeigt: \( F = -\nabla Q \).

\subsection{Die Physik ist vereinheitlicht}

Alle Theorien emergieren aus Q.

\begin{itemize}
    \item DBT = Q
    \item WG = \( -\nabla Q \)
    \item WED = \( -\nabla Q \)
\end{itemize}

\section{Die Konsequenz}

Die IWT ist die fundamentale Theorie.

Q ist die Quelle.

Alles emergiert aus Q.

Die Physik ist vereinheitlicht.

% ============================================================================
% TEIL III: DIE FOLGEN
% ============================================================================

\part{Die Folgen}

\chapter{Die Emergenz der Gravitation aus Q}
\label{ch:gravitation}

% ============================================================================
% kapitel/gravitation.tex
% ============================================================================

\section{Die Emergenz der Gravitation aus Q}

Die Gravitation ist nicht fundamental.

Sie emergiert aus Q.

\section{Die Herleitung}

Aus Q folgt:

\[
F_{\text{WG}} = -\nabla Q
\]

Die Weber-Gravitation ist der Gradient von Q.

\section{Die Bestätigung}

Die Simulation zeigt:

\[
\langle Q \rangle < 0
\]

Q wirkt bindend.

Das ist die Gravitation.

\section{Die Konsequenz}

Die Gravitation ist keine fundamentale Kraft.

Sie ist eine emergente Eigenschaft der Informationsorganisation.

Q ist die Quelle der Gravitation.

Das Graviton ist das Neutrino.

\chapter{Die Emergenz der Quantenmechanik aus Q}
\label{ch:quantenmechanik}

% ============================================================================
% kapitel/quantenmechanik.tex
% ============================================================================

\section{Die Emergenz der Quantenmechanik aus Q}

Die Quantenmechanik ist nicht fundamental.

Sie emergiert aus Q.

\section{Die Herleitung}

Aus Q folgt:

\[
Q = -\frac{\hbar^2}{2m} \frac{\nabla^2 \sqrt{\rho}}{\sqrt{\rho}}
\]

Das ist das Bohm-Potential.

Die DBT ist Q.

\section{Die Konsequenz}

Die Quantenmechanik ist eine emergente Theorie.

Sie ist eine Näherung der IWT.

Die Unvollständigkeit der QM kommt daher, dass sie die Fluktuationen nicht enthält.

\chapter{Die Emergenz der Elektrodynamik aus Q}
\label{ch:elektrodynamik}

% ============================================================================
% kapitel/elektrodynamik.tex
% ============================================================================

\section{Die Emergenz der Elektrodynamik aus Q}

Die Elektrodynamik ist nicht fundamental.

Sie emergiert aus Q.

\section{Die Herleitung}

Aus Q folgt:

\[
F_{\text{WED}} = -\nabla Q
\]

Die Weber-Elektrodynamik ist der Gradient von Q.

\section{Die Konsequenz}

Die Elektrodynamik ist eine emergente Theorie.

Sie ist eine Näherung der IWT.

Die Maxwell-Theorie ist eine Näherung der WED.

\chapter{Die Einheit der Physik}
\label{ch:einheit}

% ============================================================================
% kapitel/einheit.tex
% ============================================================================

\section{Die Einheit der Physik}

Die gesamte Physik ist Q.

\section{Die drei Säulen}

\begin{itemize}
    \item \textbf{DBT:} Q (Quantenpotential)
    \item \textbf{WG:} \( -\nabla Q \) (Gravitation)
    \item \textbf{WED:} \( -\nabla Q \) (Elektrodynamik)
\end{itemize}

\section{Die Quelle}

Q ist die Quelle.

Q ist die Organisation der Information.

Q ist das Prinzip der Selbstorganisation.

Q ist der Ursprung von allem.

\section{Die Konsequenz}

Die Physik ist vereinheitlicht.

Es gibt nur eine Theorie: die IWT.

Alles emergiert aus Q.

\section{Die Wahrheit}

Die Wahrheit ist einfach.

Information organisiert sich selbst.

Q ist der Name dieser Organisation.

Alles andere ist eine Erscheinungsform von Q.

Das ist die IWT.

% ============================================================================
% ANHANG
% ============================================================================

\appendix

\chapter{Herleitungen}
\label{app:herleitungen}

% ============================================================================
% kapitel/herleitungen.tex
% ============================================================================

\section{Herleitungen}

\subsection{Herleitung von Q}

Aus den Prinzipien der Informationserhaltung und minimaler Spannung folgt Q.

\subsection{Herleitung der DBT aus Q}

Aus Q folgt das Bohm-Potential.

\subsection{Herleitung der WG aus Q}

Aus Q folgt die Weber-Gravitation.

\subsection{Herleitung der WED aus Q}

Aus Q folgt die Weber-Elektrodynamik.

\section{Die Details}

Die vollständigen Herleitungen sind in den Anhängen der Originalarbeit enthalten.

\chapter{Numerische Details}
\label{app:numerik}

% ============================================================================
% kapitel/numerik.tex
% ============================================================================

\section{Numerische Details}

\subsection{Die Simulation}

Die Simulation verwendet:

\begin{itemize}
    \item \( N = 4096 \) Knoten
    \item \( 500 \) Iterationen
    \item GPU-Beschleunigung (OpenCL)
    \item Leapfrog-Integrator
\end{itemize}

\subsection{Die Parameter}

\begin{align*}
    \text{SCALE} &= 0.70710678 \\
    \text{ALPHA\_0} &= 1e-2 \\
    \text{RHO\_0} &= 1e-6 \\
    \text{RHO\_MIN} &= 1e-8
\end{align*}

\subsection{Die Ergebnisse}

Die Simulation zeigt:

\[
\sum |I|^2 = \text{const}
\]

\[
\langle Q \rangle < 0
\]

\[
F = -\nabla Q
\]

\chapter{Die eingefrorenen Funktionen}
\label{app:frozen}

% ============================================================================
% kapitel/frozen.tex
% ============================================================================

\section{Die eingefrorenen Funktionen}

Die Simulation hat zwei eingefrorene Funktionen.

\subsection{Fluktuationen (Vakuum)}

Die Fluktuationen wirken nur im Vakuum.

Sie sind die Quelle der Strukturbildung.

\subsection{Energiesenke (Strukturen)}

Die Energiesenke wirkt nur in Strukturen.

Sie entfernt Energie aus den Strukturen.

\subsection{Die Balance}

Die Balance zwischen Fluktuation und Energiesenke ist gefunden.

\[
\text{Fluktuation} = \text{Energiesenke}
\]

Das System ist stabil.

\subsection{Die Bedeutung}

Die eingefrorenen Funktionen sind das Herz der Simulation.

Sie zeigen, dass die IWT funktioniert.

Sie sind der Beweis für die Theorie.

% ============================================================================
% LITERATUR
% ============================================================================

\bibliographystyle{plain}
\bibliography{literatur}

\end{document}